\begin{proof}[Proof of \cref{obj:lem:transported-cace-identification}]
Fix \(P\in\mathcal M_n\) as in the statement. Throughout, conditional means are understood through their defining integral identities over measurable covariate sets.

\begin{enumerate}
\item Choose common conditional-mean versions for the transported quantities. By \cref{obj:ass:transport-complier-share,obj:ass:transport-complier-outcome}, there are measurable functions \(q_D,q_Y:\mathbb R\to\mathbb R\) satisfying
\[
\begin{aligned}
q_D(x)
&=\mathbb E_{P_n^F}[C\mid X=x,S=0]
 =\mathbb E_{P_n^F}[C\mid X=x,S=1],\\
q_Y(x)
&=\mathbb E_{P_n^F}[(Y(1)-Y(0))C\mid X=x,S=0]\\
&=\mathbb E_{P_n^F}[(Y(1)-Y(0))C\mid X=x,S=1],
\end{aligned}
\]
as versions in the respective populations.

By \cref{obj:ass:monotonicity}, the binary receipt pair
\((D(0),D(1))\) belongs almost surely to
\(\{(0,0),(0,1),(1,1)\}\). For these three pairs, respectively, both the receipt difference and the complier indicator equal \(0,1,0\). Consequently,
\[
D(1)-D(0)=C
\qquad P_n^F\text{-almost surely}.
\]
Checking the four binary receipt pairs also gives the pointwise identity
\[
Y(D(1))-Y(D(0))
=(Y(1)-Y(0))(D(1)-D(0)).
\]
Substitution therefore yields
\[
Y(D(1))-Y(D(0))=(Y(1)-Y(0))C
\qquad P_n^F\text{-almost surely}.
\]
Each population law obtained by conditioning on \(S\) is absolutely continuous with respect to \(P_n^F\). Thus these identities persist in the source population and in every restricted integral defining a conditional mean. In particular,
\[
\begin{aligned}
q_D(x)&=\mathbb E_{P_n^F}[D(1)-D(0)\mid X=x,S=1],\\
q_Y(x)&=\mathbb E_{P_n^F}[Y(D(1))-Y(D(0))\mid X=x,S=1]
\end{aligned}
\]
are source conditional-mean versions.
% lean: CausalSmith.Stat.TransportCaceRoughnuisanceLength.conditionalMean_congr_full_ae

\item Construct source arm-mean versions. Define the covariate domain, potential arm responses, zero extensions, and assignment weights by
\[
\mathcal X=[0,1],\qquad
A(z)=
\begin{cases}
D(z),&A=D,\\
Y(D(z)),&A=Y,
\end{cases}
\quad A\in\{D,Y\},\ z\in\{0,1\},
\]
\[
\widetilde m_{Az}(x)=
\begin{cases}
m_{Az}(x),&x\in\mathcal X,\\
0,&x\notin\mathcal X,
\end{cases}
\qquad
\pi_z(x)=
\begin{cases}
e(x),&z=1,\\
1-e(x),&z=0,
\end{cases}
\quad x\in\mathbb R.
\]
Also write \(P_{\mathrm S}^{X}\) for the source covariate marginal:
\[
P_{\mathrm S}^{X}(B)=P_{\mathrm S}(X\in B),
\qquad B\subseteq\mathbb R\text{ Borel}.
\]
The source marginal identification supplied by conditional randomization gives
\[
P_{\mathrm S}^{X}(B)=P_n^F(X\in B\mid S=1).
\]
By \cref{obj:ass:source-density-bounds}, this marginal is supported on \(\mathcal X\), so
\[
\widetilde m_{Az}=m_{Az}
\qquad P_{\mathrm S}^{X}\text{-almost everywhere}.
\]
The regularity and range conditions in \cref{obj:ass:arm-means-holder} make the zero extensions measurable and give
\[
0\le \widetilde m_{Az}(x)\le1
\quad(x\in\mathbb R).
\]
The binary receipts and potential-outcome bounds in
\cref{obj:ass:source-density-bounds} give, in the source population,
\[
0\le A(z)\le1
\qquad\text{almost surely}.
\]
Moreover, \cref{obj:ass:propensity-overlap} gives
\[
\frac14\le\pi_z(x)\le\frac34
\qquad(x\in\mathcal X).
\]
These bounds ensure the integrability of the variables used below.

For a Borel set \(B\subseteq\mathbb R\), put
\[
B_{\mathcal X}=B\cap\mathcal X.
\]
The source arm-moment identity in \cref{obj:ass:arm-means-holder} gives
\[
\int_{B_{\mathcal X}}
 f_{\mathrm S}(x)\pi_z(x)m_{Az}(x)\,dx
=
\mathbb E_{P_{\mathrm S}}
 \bigl[\mathbf1_{\{X\in B_{\mathcal X},\,Z=z\}}A\bigr].
\]
Receipt and outcome consistency, as supplied by
\cref{obj:ass:receipt-consistency,obj:ass:outcome-consistency},
replace the observed response on the arm \(Z=z\) by \(A(z)\).
Conditional randomization in \cref{obj:ass:instrument-randomization}
makes the full-data law restricted to that arm have density
\(\pi_z(X)\) relative to the source full-data law. Hence
\[
\begin{aligned}
\mathbb E_{P_{\mathrm S}}
 \bigl[\mathbf1_{\{X\in B_{\mathcal X},\,Z=z\}}A\bigr]
&=\mathbb E
 \bigl[\mathbf1_{\{X\in B_{\mathcal X},\,Z=z\}}A(z)\bigr]\\
&=\mathbb E_{P_n^F}
 \bigl[\mathbf1_{\{X\in B_{\mathcal X}\}}\pi_z(X)A(z)
       \mid S=1\bigr],
\end{aligned}
\]
where the expectation in the middle is under the source assignment law. The last equality follows first for indicators of measurable full-data events from the assignment probability identity, and then for the displayed bounded measurable integrand by integration.

Using the source density identity and source support, the preceding equalities become
\[
\begin{aligned}
&\mathbb E_{P_n^F}
 \bigl[\mathbf1_{\{X\in B\}}\pi_z(X)\widetilde m_{Az}(X)
       \mid S=1\bigr]\\
&\qquad=
\int_{B_{\mathcal X}}
 f_{\mathrm S}(x)\pi_z(x)m_{Az}(x)\,dx\\
&\qquad=
\mathbb E_{P_n^F}
 \bigl[\mathbf1_{\{X\in B\}}\pi_z(X)A(z)\mid S=1\bigr].
\end{aligned}
\]
Thus \(\pi_z(X)\widetilde m_{Az}(X)\) is a version of the conditional mean of \(\pi_z(X)A(z)\) given \(X\) in the source population. Since \(\pi_z(X)\) is measurable with respect to \(X\), the conditional-expectation multiplication rule gives
\[
\pi_z(X)\widetilde m_{Az}(X)
=
\mathbb E_{P_n^F}[\pi_z(X)A(z)\mid X,S=1]
=
\pi_z(X)\mathbb E_{P_n^F}[A(z)\mid X,S=1].
\]
The positive lower bound on \(\pi_z(X)\) permits cancellation:
\[
\widetilde m_{Az}(X)
=\mathbb E_{P_n^F}[A(z)\mid X,S=1]
\qquad\text{almost surely in the source population}.
\]
Equivalently, for every Borel \(B\subseteq\mathbb R\),
\[
\mathbb E_{P_n^F}
 \bigl[\mathbf1_{\{X\in B\}}A(z)\mid S=1\bigr]
=
\int_B\widetilde m_{Az}(x)\,dP_{\mathrm S}^{X}(x).
\]
% lean: CausalSmith.Stat.TransportCaceRoughnuisanceLength.armMeanVersion_conditionalMean

\item Identify the contrasts with the common versions. Define, for \(A\in\{D,Y\}\),
\[
\widetilde\Delta_A(x)
=\widetilde m_{A1}(x)-\widetilde m_{A0}(x),
\qquad
q_A=
\begin{cases}
q_D,&A=D,\\
q_Y,&A=Y.
\end{cases}
\]
The source almost-everywhere equalities for the arm means imply
\[
\Delta_A=\widetilde\Delta_A
\qquad P_{\mathrm S}^{X}\text{-almost everywhere}.
\]
Subtracting the two arm integral identities from the preceding step gives
\[
\int_B\widetilde\Delta_A(x)\,dP_{\mathrm S}^{X}(x)
=
\mathbb E_{P_n^F}
 \bigl[\mathbf1_{\{X\in B\}}(A(1)-A(0))\mid S=1\bigr]
=
\int_Bq_A(x)\,dP_{\mathrm S}^{X}(x)
\]
for every Borel \(B\subseteq\mathbb R\).

To obtain almost-everywhere equality from these identities, define
\[
\mathcal E_{k,A}
=
\bigl\{x\in\mathbb R:
 |\widetilde\Delta_A(x)|\le k,\ |q_A(x)|\le k\bigr\},
\qquad k\in\mathbb N.
\]
Both functions are measurable and real-valued, and
\[
\bigcup_{k\in\mathbb N}\mathcal E_{k,A}=\mathbb R.
\]
On each \(\mathcal E_{k,A}\), they are integrable under the finite source covariate law. Applying the preceding integral equality to
\(B\cap\mathcal E_{k,A}\), uniqueness for integrable functions with equal integrals over every measurable set yields
\[
\widetilde\Delta_A=q_A
\qquad
P_{\mathrm S}^{X}\text{-almost everywhere on }\mathcal E_{k,A}.
\]
Taking the countable union, and then using the equality between the raw and extended contrasts, gives
\[
\Delta_D=q_D,\qquad \Delta_Y=q_Y
\qquad P_{\mathrm S}^{X}\text{-almost everywhere}.
\]

The density identities in
\cref{obj:ass:source-density-bounds,obj:ass:target-density-bounds}
transfer these equalities to the target covariate law. Indeed, if a Borel set \(B\) is source-null, then
\[
0=P_{\mathrm S}^{X}(B)
=\int_{B\cap\mathcal X}f_{\mathrm S}(x)\,dx
\ge c_f\int_{B\cap\mathcal X}1\,dx.
\]
Since \(c_f>0\), the set \(B\cap\mathcal X\) is Lebesgue-null, and consequently
\[
P_{\mathrm T}(B)
=\int_{B\cap\mathcal X}f_{\mathrm T}(x)\,dx=0.
\]
Thus
\[
P_{\mathrm T}\ll P_{\mathrm S}^{X},
\qquad
\Delta_D=q_D,\quad\Delta_Y=q_Y
\quad P_{\mathrm T}\text{-almost everywhere}.
\]
Together with the conditional-mean identities for \(q_D,q_Y\), this supplies common versions for both chains in the statement.
% lean: CausalSmith.Stat.TransportCaceRoughnuisanceLength.target_absolutelyContinuous_source

\item Integrate the common versions in the target population. The target density identity and its positive envelope in
\cref{obj:ass:target-density-bounds} convert the defining transported integrals into integrals under \(P_{\mathrm T}\). Applying the target conditional-mean identities to the whole covariate space gives
\[
\begin{aligned}
T_D
&=\int_{\mathcal X}f_{\mathrm T}(x)\Delta_D(x)\,dx\\
&=\int_{\mathbb R}\Delta_D(x)\,dP_{\mathrm T}(x)
 =\int_{\mathbb R}q_D(x)\,dP_{\mathrm T}(x)\\
&=\mathbb E_{P_n^F}[C\mid S=0]
 =P_n^F(C=1\mid S=0)=\mu.
\end{aligned}
\]
The penultimate equality uses
\(C=\mathbf1_{\{C=1\}}\), so its integral is the probability of the complier event. Similarly,
\[
\begin{aligned}
T_Y
&=\int_{\mathcal X}f_{\mathrm T}(x)\Delta_Y(x)\,dx\\
&=\int_{\mathbb R}\Delta_Y(x)\,dP_{\mathrm T}(x)
 =\int_{\mathbb R}q_Y(x)\,dP_{\mathrm T}(x)\\
&=\mathbb E_{P_n^F}[(Y(1)-Y(0))C\mid S=0].
\end{aligned}
\]
% lean: CausalSmith.Stat.TransportCaceRoughnuisanceLength.transported_cace_identification

\item Complete the arm identities and the effect-domain assertion. For each \(z\in\{0,1\}\), the constructed versions satisfy
\[
\begin{aligned}
m_{Dz}(x)
&=\widetilde m_{Dz}(x)
 =\mathbb E_{P_n^F}[D(z)\mid X=x,S=1],\\
m_{Yz}(x)
&=\widetilde m_{Yz}(x)
 =\mathbb E_{P_n^F}[Y(D(z))\mid X=x,S=1]
\end{aligned}
\]
for \(P_{\mathrm S}^{X}\)-almost every \(x\). These are the required source arm identities.

By the definition of the target complier effect and the two transported integral identities,
\[
\theta
=
\frac{\mathbb E_{P_n^F}[(Y(1)-Y(0))C\mid S=0]}
     {P_n^F(C=1\mid S=0)}
=
\frac{T_Y}{T_D}.
\]
By \cref{obj:ass:positive-strength}, \(T_D=\mu>0\).

Finally, the potential-outcome support restriction in
\cref{obj:ass:source-density-bounds} persists under the target population law. On the event \(C=1\),
\[
-1\le Y(1)-Y(0)\le1,
\]
whereas on \(C=0\) the complier-weighted effect is zero. Hence, target-almost surely,
\[
-C\le (Y(1)-Y(0))C\le C.
\]
All three variables are integrable by their bounds. Integration therefore gives
\[
-\mu
\le\mathbb E_{P_n^F}[(Y(1)-Y(0))C\mid S=0]
=T_Y
\le\mu.
\]
Dividing by \(\mu=T_D>0\) yields
\[
-1\le\frac{T_Y}{T_D}=\theta\le1,
\qquad\text{so}\qquad
\theta\in[-1,1]=\Theta.
\]
% lean: CausalSmith.Stat.TransportCaceRoughnuisanceLength.targetCACE_mem_parameterSpace
\end{enumerate}
\end{proof}
